\documentclass[11pt,a4paper]{article}

\usepackage[utf8]{inputenc}
\usepackage[T1]{fontenc}
\usepackage[spanish]{babel}
\usepackage[margin=2cm, top=2.6cm, bottom=1.8cm, footskip=0.6cm]{geometry}
\usepackage{graphicx}
\graphicspath{{../../01_Propuesta/Logos/}{./}}
\usepackage{fancyhdr}
\usepackage{xcolor}
\usepackage[hidelinks]{hyperref}
\usepackage{enumitem}
\setlist{nosep}
\setlength{\parindent}{0pt}
\setlength{\parskip}{4pt}

\definecolor{verdeosc}{RGB}{26,107,60}
\definecolor{verdemid}{RGB}{46,139,87}
% Datos que faltan: se ven en rojo hasta completarlos.
\newcommand{\pend}[1]{\textcolor{red}{[PENDIENTE: #1]}}

\setlength{\headheight}{40pt}
\pagestyle{fancy}
\fancyhf{}
\fancyhead[L]{%
  \begin{tabular*}{\textwidth}{@{\extracolsep{\fill}}cccc@{}}
    \includegraphics[height=1.0cm]{logo_unisucre.png} &
    \includegraphics[height=1.0cm]{logo_sicvec2026.png} &
    \includegraphics[height=1.0cm]{Logo_insilico.jpeg} &
    \includegraphics[height=1.0cm]{Logo_Bio.png} \\
  \end{tabular*}%
}
\fancyfoot[C]{\footnotesize\itshape SICVEC 2026 contribuye a los Objetivos de Desarrollo Sostenible 12, 13, 15 y 17.}
\renewcommand{\headrulewidth}{1pt}
\renewcommand{\headrule}{\hbox to\headwidth{\color{verdeosc}\leaders\hrule height\headrulewidth\hfill}}

\begin{document}

\vspace{0.1cm}

\noindent
\textbf{Sincelejo, 29 de septiembre de 2026}

\vspace{0.1cm}

\noindent
Señor\\
Esp. Carlos Meza Bertel\\
Director Ambiental y de Servicios Generales\\
Parque Comercial Guacarí\\
Sincelejo, Sucre, Colombia

\vspace{0.1cm}

\noindent
\textbf{\large Asunto: Invitación como ponente — SICVEC 2026}

\vspace{0.1cm}

Estimado señor Meza:

\vspace{0.1cm}

En nombre del Grupo de Investigación IN SILICO y del Departamento de Biología y Química de la Universidad de Sucre (UNISUCRE), me complace invitarlo a participar como \textbf{ponente} en el \textbf{Simposio Internacional de Ciencia Verde y Economía Circular (SICVEC 2026)}, que se realizará en el propio Centro Comercial Guacarí.

Desde la Dirección Ambiental y de Servicios Generales del Parque Comercial Guacarí, usted conoce de primera mano cómo se gestionan los residuos, el agua, la energía y los servicios de un gran espacio comercial en Sincelejo. Esa experiencia es la que queremos llevar al simposio: una mirada práctica a la economía circular y la sostenibilidad desde la operación diaria de una empresa de la región, que complemente las ponencias académicas y muestre a estudiantes, investigadores y empresarios cómo se aplican estos principios fuera del laboratorio.

\vspace{0.1cm}

\noindent
\textbf{INFORMACIÓN DEL EVENTO}

\begin{tabular}{ll}
\textbf{Nombre:} & Simposio Internacional de Ciencia Verde y Economía Circular \\
\textbf{Fechas:} & 19 y 20 de octubre de 2026 (lunes y martes) \\
\textbf{Lugar:} & Centro Comercial Guacarí, Sincelejo, Sucre, Colombia \\
\textbf{Modalidad:} & Híbrida \\
\textbf{Sitio web:} & \href{https://sicvec2026.eu.pythonanywhere.com}{sicvec2026.eu.pythonanywhere.com} \\
\end{tabular}

\vspace{0.1cm}

\noindent
\textbf{DETALLES DE LA PONENCIA}

\begin{tabular}{ll}
\textbf{Duración:} & 20 minutos \\
\textbf{Eje temático:} & Economía Circular \\
\textbf{Tema sugerido:} & Economía circular y sostenibilidad desde el Parque Comercial Guacarí \\
\textbf{Participación:} & Presencial, en el Centro Comercial Guacarí \\
\textbf{Franja:} & Martes 20 de octubre, 08:30--09:00 \\
\textbf{Modalidad:} & Invitación académica (sin honorarios) \\
\end{tabular}

\vspace{0.1cm}

\noindent
\textbf{PRÓXIMOS PASOS}

\begin{enumerate}
  \item Confirmar su disponibilidad
  \item Proponer título y un breve resumen de su charla
  \item Enviar una reseña corta y una foto para la promoción del evento
\end{enumerate}

Rogamos responder esta invitación \textbf{lo antes posible} para incluir su charla en el programa.

\vspace{0.1cm}

Quedo atento a sus comentarios y a su disposición para coordinar cualquier detalle. Será un gusto contar con la experiencia del Parque Comercial Guacarí en SICVEC 2026.

\vspace{0.15cm}

Atentamente,

\vspace{0.15cm}

\textbf{Dr. Aldo F. Combariza}\\
Coordinador General, SICVEC 2026\\
Grupo de Investigación IN SILICO · Departamento de Biología y Química\\
Facultad de Educación y Ciencias · Universidad de Sucre\\
\href{mailto:aldo.combariza@unisucre.edu.co}{aldo.combariza@unisucre.edu.co}

\vspace{0.1cm}

\end{document}
